\documentclass[10pt,a4paper]{article}
\usepackage[margin=20mm]{geometry}
\usepackage{fontspec}
\usepackage{amsmath,amssymb}
\usepackage{booktabs}
\usepackage{tabularx}
\usepackage{enumitem}
\usepackage{xcolor}
\usepackage{hyperref}
\usepackage{listings}

\setmainfont{Noto Sans CJK KR}[AutoFakeSlant=0.2]
\setsansfont{Noto Sans CJK KR}
\setmonofont{D2Coding}
\hypersetup{
  colorlinks=true,
  linkcolor=blue,
  urlcolor=blue,
  pdftitle={3번 네트워크보안: TLS 1.2 AES-GCM Nonce 재사용을 이용한 Record 위조},
  pdfauthor={암호분석경진대회 제출 보고서},
  pdfsubject={TLS 1.2 AES-GCM nonce reuse and GHASH forgery},
  pdfkeywords={TLS 1.2, AES-GCM, nonce reuse, GHASH, forgery}
}
\setlength{\parindent}{0pt}
\setlength{\parskip}{5pt}
\setlist{nosep,leftmargin=6mm}
\lstset{
  basicstyle=\ttfamily\small,
  frame=single,
  columns=fullflexible,
  breaklines=true,
  keepspaces=true
}

\title{3번 네트워크보안\\TLS 1.2 AES-GCM Nonce 재사용을 이용한 Record 위조}
\author{암호분석경진대회 제출 보고서}
\date{2026년 7월 28일}

\begin{document}
\maketitle

\section{결과 요약}

제공된 \texttt{tls\_live.pcap}의 서버 방향 TLS 1.2 통신은
\texttt{TLS\_ECDHE\_RSA\_WITH\_AES\_128\_GCM\_SHA256}을 사용한다.
서버 application data 1--3번이 같은 8바이트 explicit nonce를
재사용하므로, 같은 write key와 4바이트 fixed salt 아래에서 전체
96비트 GCM nonce가 충돌한다. 이 충돌로 CTR keystream을 재사용해
\texttt{amt=0100}을 \texttt{amt=0500}으로 바꾸고, tag 차분에서
GHASH subkey를 복원해 유효한 새 tag를 계산했다.

최종 제출 record는 77바이트이다. 아래 네 줄을 개행 없이 연결한 값을 제출한다.

\begin{lstlisting}
17030300489c269a9f29810ab0
a99141d1d84c3df6bfdcdb7dec7cf30462db203fbc36a29772edbf159643b468
f1f1e520e7341daee867a49665542cd9
f01f5a9c906b73ba3e13bbfea0e98a18
\end{lstlisting}

\section{PCAP 분석과 취약점 식별}

\subsection{TCP/TLS record 복원}

pcap의 link type은 \texttt{DLT\_IPV4(228)}이다. 서버 방향 flow는
다음과 같다.

\begin{center}
\texttt{10.0.0.1:443 -> 10.0.0.2:54321}
\end{center}

TCP segment를 sequence number 순으로 정렬하고 overlap을 제거해 byte stream을 재조립했다.
5바이트 TLS header의 길이 필드로 record를 분리하고 ServerHello에서
cipher suite \texttt{0xc02f}를 확인했다.

ChangeCipherSpec 뒤 새 server write state의 TLS sequence number는
0으로 초기화된다. encrypted Finished가 \texttt{seq=0}이고, 이어지는
application data의 sequence number는 1--4이다.

\begin{center}
\small
\begin{tabular}{rrrrl}
\toprule
App record & TLS seq & Ciphertext & Fragment & Explicit nonce \\
\midrule
1 & 1 & 34 & 58 & \texttt{9c269a9f29810ab0} \\
2 & 2 & 50 & 74 & \texttt{9c269a9f29810ab0} \\
3 (target) & 3 & 48 & 72 & \texttt{9c269a9f29810ab0} \\
4 & 4 & 40 & 64 & \texttt{9c269a9f29810ab1} \\
\bottomrule
\end{tabular}
\end{center}

\subsection{TLS 1.2 GCM nonce와 AAD}

TLS 1.2 AES-GCM nonce는 write state의 4바이트 fixed salt와 record의
8바이트 explicit nonce를 연결한 값이다.

\[
  \mathit{nonce}=\mathit{write\_salt}[4]\parallel
  \mathit{nonce\_explicit}[8].
\]

세 record는 같은 연결, 방향, cipher epoch에 있으므로 key와 salt가
같다. 따라서 같은 explicit nonce는 같은 전체 nonce, 같은 CTR
keystream, 같은 $S=E_K(J_0)$를 뜻한다.

TLS AEAD additional data는 다음 13바이트다.

\[
  A=\mathit{seq}[8]\parallel\mathit{type}[1]\parallel
  \mathit{version}[2]\parallel\mathit{plaintext\_length}[2].
\]

목표 record의 AAD는 다음과 같다.

\begin{lstlisting}
0000000000000003 || 17 || 0303 || 0030
= 00000000000000031703030030
\end{lstlisting}

AAD의 \texttt{0030}은 48바이트 평문 길이이다. record header의
\texttt{0048}은 explicit nonce 8바이트, ciphertext 48바이트, tag
16바이트를 합한 72바이트 wire fragment 길이이므로 서로 다르다.
GCM에는 CBC식 TLS padding이나 별도 HMAC이 없다. GHASH 계산 과정의
zero padding만 존재하며 wire에는 전송되지 않는다.

\section{GHASH Authentication Subkey 복원}

\subsection{Nonce 재사용 방정식}

GCM에서

\[
  H=E_K(0^{128}),\qquad
  T_i=S\oplus\operatorname{GHASH}_H(A_i,C_i)
\]

이다. 같은 nonce인 record 두 개의 tag를 XOR하면 $S$가 제거된다.

\[
 T_1\oplus T_2=
 \operatorname{GHASH}_H(A_1,C_1)\oplus
 \operatorname{GHASH}_H(A_2,C_2).
\]

GHASH는

\[
\mathbb F_{2^{128}}=
\mathbb F_2[x]/(x^{128}+x^7+x^2+x+1)
\]

에서 계산한다. block을 $X_1,\ldots,X_m$이라 하면

\[
Y_0=0,\qquad Y_i=(Y_{i-1}\oplus X_i)H
\]

이므로 각 GHASH는 $H$에 대한 다항식 $P_i(H)$이다. 2번 ciphertext는
50바이트라서 AAD 1 block, ciphertext 4 blocks, length 1 block을
사용한다. 따라서 1번과 2번 record로 만든

\[
F_{12}(H)=P_1(H)\oplus P_2(H)\oplus T_1\oplus T_2=0
\]

은 6차 방정식이다.

\subsection{유한체 근 분리와 검증}

solver는

\[
\gcd\!\left(F_{12}(X),X^{2^{128}}-X\right)
\]

로 $\mathbb F_{2^{128}}$ 안의 선형 인수만 분리한 뒤 근을 나눈다.
NIST GCM은 field element를 MSB-first로 표현하고 구현의 일반 다항식
연산은 bit 0을 상수항 계수로 사용하므로, 표현 경계에서 128비트
bit reverse를 적용했다. 구체적으로 NIST 표현의 왼쪽 bit가 $x^0$의
계수이므로 다음 경계 대응을 사용한다.

\begin{lstlisting}
NIST 0x80000000000000000000000000000000 <-> polynomial integer 1
NIST 0x40000000000000000000000000000000 <-> polynomial integer 2 (= x)
\end{lstlisting}

solver는 첫 번째 값을 GCM 곱셈 항등원으로 사용해 임의 probe와 곱한
결과가 probe 그대로인지 실행 시 self-test한다. 1번과 2번 record로 만든
6차식에서 체 내부 root는 1개였고, 3번 record의 동일 $S$ 조건까지
통과한 후보도 1개였다.

\begin{center}
\begin{tabular}{lr}
\toprule
검증 단계 & 후보 수 \\
\midrule
$F_{12}$의 $\mathbb F_{2^{128}}$ root & 1 \\
3번 record의 동일 $S$ 조건 통과 & 1 \\
\bottomrule
\end{tabular}
\end{center}

각 후보는 재사용 record 세 개 모두에서

\[
T_i\oplus P_i(H)=S
\]

가 같은지 검사했다. 유일하게 남은 값은 다음과 같다.

\begin{lstlisting}
H        = 112332a84132bc0c5c23a61037723683
E_K(J0)  = 1e023d8461003896981f1441c5feab90
\end{lstlisting}

$H$는 AES session key가 아니라 GCM authentication subkey이다. 하지만
재사용 nonce의 $H$와 $S$를 모두 얻었으므로 그 nonce에 대한 새
ciphertext와 tag를 만들 수 있다.

\section{평문 및 Tag 위조}

\subsection{Ciphertext 변경}

원래 평문은 문제 본문에 직접 주어진 known plaintext이며, pcap을
별도로 복호화해 얻은 값이 아니다. 두 평문은 모두 48바이트이고
0-based offset 32의 한 바이트만 다르다.

\begin{lstlisting}
old: action=set_salary&uid=0007&amt=0100&month=202603
new: action=set_salary&uid=0007&amt=0500&month=202603
\end{lstlisting}

\texttt{'1'=0x31}, \texttt{'5'=0x35}이므로 delta는 \texttt{0x04}다.
CTR에서 $C=P\oplus KS$이므로

\[
C'=C\oplus P\oplus P'
\]

로 바꾸면 된다. ciphertext의 해당 byte는
\texttt{f5 xor 04 = f1}이 된다.

\begin{lstlisting}
old C = a99141d1d84c3df6bfdcdb7dec7cf30462db203fbc36a29772edbf159643b468f5f1e520e7341daee867a49665542cd9
new C = a99141d1d84c3df6bfdcdb7dec7cf30462db203fbc36a29772edbf159643b468f1f1e520e7341daee867a49665542cd9
\end{lstlisting}

\subsection{Tag 재계산}

nonce, AAD, 길이가 그대로이므로 $S$도 같다.

\[
\begin{aligned}
S&=T_{\rm old}\oplus\operatorname{GHASH}_H(A,C_{\rm old}),\\
T_{\rm new}&=S\oplus\operatorname{GHASH}_H(A,C_{\rm new}).
\end{aligned}
\]

\begin{lstlisting}
old tag   = b0d6e0ac2be205621f7ea4a034519f95
GHASH old = aed4dd284ae23df48761b0e1f1af3405
S         = 1e023d8461003896981f1441c5feab90
GHASH new = ee1d6718f16b4b2ca60cafbf65172188
new tag   = f01f5a9c906b73ba3e13bbfea0e98a18
\end{lstlisting}

\section{최종 Record 구조}

\begin{center}
\small
\begin{tabularx}{\textwidth}{rrlX}
\toprule
Offset & Size & Field & Value \\
\midrule
0 & 1 & ContentType & \texttt{17} \\
1--2 & 2 & TLS version & \texttt{0303} \\
3--4 & 2 & Fragment length & \texttt{0048} \\
5--12 & 8 & Explicit nonce & \texttt{9c269a9f29810ab0} \\
13--60 & 48 & Forged ciphertext & \texttt{a99141...5542cd9} \\
61--76 & 16 & GCM tag & \texttt{f01f5a9c906b73ba3e13bbfea0e98a18} \\
\bottomrule
\end{tabularx}
\end{center}

fragment는 $8+48+16=72$바이트이고 header를 포함한 record는
$5+72=77$바이트, 즉 개행을 제외한 hex 154문자다.

\section{구현과 재현 검증}

구현 파일은 \texttt{solve\_03\_tls\_gcm\_nonce\_reuse.py}다. Python 표준
라이브러리만 사용하며 pcap magic/endian/link type과 packet 길이를 검사하고,
TCP segment를 sequence number로 재조립하며, TLS record 경계를
검증한다. GF($2^{128}$) 곱셈, 역원, 다항식 GCD와 인수분해도 직접
구현했다. 인수분해 난수 seed는 3으로 고정하고 최종 후보는 세 tag
방정식으로 결정론적으로 검증한다.

저장소 루트에서 다음을 실행한다.

\begin{lstlisting}
python3 solutions/solve_03_tls_gcm_nonce_reuse.py
\end{lstlisting}

코드는 다음 조건을 모두 assertion으로 검사한다.

\begin{itemize}
  \item cipher suite, TLS sequence number, nonce 재사용 pattern;
  \item 유일한 $H$와 세 record에서 동일한 $S$;
  \item 새 ciphertext가 재사용 keystream 아래 목표 평문으로 복원됨;
  \item $T_{\rm new}\oplus\operatorname{GHASH}_H(A,C_{\rm new})=S$;
  \item fragment/record 길이와 최종 예상 hex; 그리고
  \item 생성값과 \texttt{submissions/03/forged\_record.txt}의 일치.
\end{itemize}

\section{참고 자료}

\begin{itemize}
  \item \href{https://www.rfc-editor.org/rfc/rfc5246.html#section-6.2.3.3}
        {IETF RFC 5246, \emph{The Transport Layer Security (TLS) Protocol
        Version 1.2}}, Section 6.2.3.3.
  \item \href{https://www.rfc-editor.org/rfc/rfc5288.html#section-3}
        {IETF RFC 5288, \emph{AES Galois Counter Mode (GCM) Cipher Suites for
        TLS}}, Section 3.
  \item \href{https://doi.org/10.6028/NIST.SP.800-38D}
        {NIST SP 800-38D, \emph{Recommendation for Block Cipher Modes of
        Operation: Galois/Counter Mode (GCM) and GMAC}}.
  \item \href{https://eprint.iacr.org/2004/193}
        {David A. McGrew and John Viega, \emph{The Security and Performance of
        the Galois/Counter Mode (GCM) of Operation}}, 2004.
  \item \href{https://csrc.nist.gov/CSRC/media/Projects/Block-Cipher-Techniques/documents/BCM/Comments/800-38-series-drafts/GCM/Joux_comments.pdf}
        {Antoine Joux, \emph{Authentication Failures in NIST version of GCM}},
        2006.
\end{itemize}

\end{document}
